\clearpage
\section{Communion Under Both Kinds and Sacramental Vessels}
Plate X unites the in-use arrangement of Communion under both kinds on an unfolded corporal with the three holy oils consecrated and blessed annually at the Chrism Mass (Sacred Chrism, Oil of Catechumens, and Oil of the Sick) and the baptismal font that stands at the threshold of Christian initiation.

\PCDPlateHeader{Plate X --- Communion Under Both Kinds and Sacramental Vessels}{Not to common scale}
\ifdefined\TriptychPrintEdition
\begin{center}
\includegraphics[width=0.92\linewidth,height=3.55in,keepaspectratio]{liturgy/roman-rite/postconciliar/2008-latin-2011-us-english/reference/roman-sanctuary-dictionary/shared/artwork/graphite/pc-art-10-communion-triduum-vessels.png}
\end{center}
\PCDPositionKey{Top: Communion under both kinds arrangement on an unfolded Corporal (principal Chalice, Paten, veiled covered Ciborium, and two footed Communion Cups each with a folded Purificator); bottom left: three cylindrical lidded metal Holy Oils Stocks on a tray (Sacred Chrism, Oil of Catechumens, Oil of the Sick); bottom right: carved octagonal stone Baptismal Font with removable domed cover ajar.}
\else
\textbf{Visual Plate Summary (Plate X --- Communion Under Both Kinds and Sacramental Vessels, Not to common scale):} Three-group monochrome graphite sacramental plate: top, an unfolded square corporal holding a central principal chalice, shallow paten, closed ciborium draped in a white veil, and two footed Communion cups each beside a folded purificator; bottom left, three matching cylindrical lidded metal holy-oil vessels resting on an oval metal tray without baked inscriptions; bottom right, an octagonal carved stone baptismal font on a pedestal with its domed cover resting slightly open to reveal the water basin rim.

\textbf{Reading order and object placement:} Top: Communion under both kinds arrangement on an unfolded Corporal (principal Chalice, Paten, veiled covered Ciborium, and two footed Communion Cups each with a folded Purificator); bottom left: three cylindrical lidded metal Holy Oils Stocks on a tray (Sacred Chrism, Oil of Catechumens, Oil of the Sick); bottom right: carved octagonal stone Baptismal Font with removable domed cover ajar.
\fi

\label{pc-obj-holy-oils-stocks}
\PCDEntryHead{Holy Oils and Oil Stocks}{Olea sancta et vasa oleorum}{required}{N40}
\PCDField{Also called}{Sacred Chrism (SC), Oil of Catechumens (OC), and Oil of the Sick (OI); oil stocks; chrismatory; oil urns; ambry (oil cabinet); \textit{Sanctum Chrisma}; \textit{Oleum catechumenorum}; \textit{Oleum infirmorum}; \textit{ampullae oleorum}; \textit{thecae}; \textit{armarium oleorum}}
The three Holy Oils---Sacred Chrism (SC), Oil of Catechumens (OC), and Oil of the Sick (OI)---are olive or other plant oils consecrated or blessed annually by the diocesan bishop for sacramental anointing, kept in matching labeled metal or glass oil stocks and safeguarded in a locked ambry (armarium oleorum). Matching cylindrical metal stocks with tightly fitted domed lids and rounded finials stand together on a rimmed tray; larger display urns hold the parish supply in the locked ambry, while smaller screw-top pocket stocks lined with oil-saturated cotton serve baptisms and sick calls. In parish use, actual stocks and lids bear unmistakable marks---SC or CHR for Sacred Chrism, OC (or older Latin OS, Oleum Sanctum---never to be confused with English ``Oil of the Sick'') for Oil of Catechumens, and OI or INF for Oil of the Sick---so that vessels and caps are never transposed.

\PCDField{Plate position and status}{Bottom left, three matching cylindrical lidded metal oil stocks on an oval tray (Plate X) (Universal Roman law)}
\PCDField{Placement and liturgical use}{At the annual Chrism Mass (Missa Chrismatis), celebrated on Holy Thursday morning or an earlier day in Holy Week close to Easter, the diocesan bishop blesses the Oil of the Sick and the Oil of Catechumens and mixes balsam or fragrance into oil to consecrate Sacred Chrism (Ordo benedicendi oleum 1970; Caeremoniale Episcoporum 1984 no. 274). Parishes receive fresh oils each year and reserve them under lock in a clean, dignified ambry in the sanctuary, baptistery, or sacristy (CIC can. 847, par. 2). Sacred Chrism anoints the newly baptized, the confirmed, newly ordained priests and bishops, and dedicated churches and altars; Oil of Catechumens is used in pre-baptismal anointing; and Oil of the Sick is administered in the Anointing of the Sick.}
\PCDField{United States norm}{Universal canon law (CIC can. 847) requires every parish to obtain newly consecrated and blessed oils each year from its diocesan bishop and safeguard them diligently. In the United States, parishes commonly receive the new oils during the Evening Mass of the Lord's Supper on Holy Thursday using the USCCB Order for the Reception of the Holy Oils, and many churches display the locked ambry visibly near the baptismal font or sanctuary.}
\PCDField{Ministers and server boundary}{Ordinary ministers: Bishop, Priest, Deacon. Server moves only at celebrant's or sacristan's direction: Carry closed oil urns in procession only when designated (such as at the Chrism Mass or Holy Thursday reception of the oils) and hold a service tray with cotton and towel for the minister during Baptism or Confirmation; never open a stock, anoint anyone, transpose caps, or handle the ambry key.}
\PCDSafety{Never open an oil stock, attempt to anoint anyone, or interchange caps between SC, OC, and OI vessels. Hold closed oil urns upright with both hands during processions to prevent oil spills on sanctuary steps, and never discard previous year's oils in ordinary sinks, trash, or plumbing; old oils are burned reverently (commonly absorbed in cotton or consumed in the Easter fire) or poured into the sacrarium or clean earth according to diocesan direction.}
\PCDField{Distinguish from}{Pyx, Cruets}
\PCDField{Symbolism}{The metal stocks themselves are practical protective containers; sacramental symbolism belongs to the consecrated and blessed oils they hold. In the Church's liturgy, oil signifies cleansing, healing, strength, and spiritual joy, while balsam-perfumed Sacred Chrism signifies the seal of the Holy Spirit conforming the anointed to Christ as priest, prophet, and king and spreading the ``fragrance of innocence'' prayed for in the Roman Chrism Preface. St. Paul writes, ``Now he that confirmeth us with you in Christ and that hath anointed us, is God: Who also hath sealed us and given the pledge of the Spirit in our hearts'' (2 Cor 1:21--22), and St. James instructs the Church concerning the sick, ``Is any man sick among you? Let him bring in the priests of the church and let them pray over him, anointing him with oil in the name of the Lord. And the prayer of faith shall save the sick man'' (James 5:14--15). (CCC 1294--1300; Ordo benedicendi oleum (1970); Roman Chrism Preface; 2 Cor 1:21--22; James 5:14--15)}
\PCDField{1962 Roman Missal counterpart}{\PCDTag{materially changed} The 1962 rites required pure olive oil exclusively, blessed strictly on Holy Thursday morning, and poured both Oil of Catechumens and Sacred Chrism directly into the Easter Vigil baptismal font. Postconciliar discipline (Ordo benedicendi oleum 1970; CIC 1983 can. 847, par. 1) permits oil pressed from olives or other plants, allows moving the Chrism Mass to an earlier day in Holy Week to foster presbyteral participation, discontinues pouring oils into the font water, and encourages prominent reservation in a locked ambry near the baptistery or sanctuary.}
\PCDEntryRule
